%-------------------------
% Two-page CV for quantum algorithms / complexity PhD applications
% Arul Rhik Mazumder
%-------------------------
\documentclass[letterpaper,10pt]{article}

\usepackage[letterpaper,margin=0.5in]{geometry}
\usepackage[T1]{fontenc}
\usepackage{titlesec}
\usepackage[usenames,dvipsnames]{color}
\usepackage{enumitem}
\usepackage[hidelinks]{hyperref}
\usepackage{graphicx}
\usepackage[english]{babel}
\usepackage{fontawesome5}

\definecolor{softblue}{RGB}{70, 110, 180}
\hypersetup{colorlinks=true, urlcolor=softblue, linkcolor=softblue, citecolor=softblue}

\pagestyle{empty}
\urlstyle{same}
\raggedbottom
\raggedright
\sloppy
\setlength{\parindent}{0pt}

\titleformat{\section}{\scshape\raggedright\large\bfseries}{}{0em}{}[\color{black}\titlerule]
\titlespacing*{\section}{0pt}{18pt}{7pt}
\titleformat{\subsection}{\raggedright\normalsize\bfseries}{}{0em}{}
\titlespacing*{\subsection}{0pt}{3pt}{1pt}

% Compact lists
\setlist[itemize]{leftmargin=1.1em, itemsep=4pt, parsep=0pt, topsep=1pt, partopsep=0pt}
\setlist[enumerate]{leftmargin=1.4em, itemsep=0.75pt, parsep=0pt, topsep=1pt, partopsep=0pt}

% Entry heading: name and dates on line 1, role and place on line 2
\newcommand{\entry}[4]{%
  \par\vspace{4pt}\noindent\textbf{#1}\hfill\textbf{\small #2}\\
  \textit{\small #3}\hfill\textit{\small #4}\par}

\begin{document}

%----------HEADER----------
\begin{center}
  {\huge\scshape Arul Rhik Mazumder}\\[1pt]
  U.S. Citizen\\[1pt]
  \small \raisebox{-0.1\height}\faPhone\ 1-978-710-1802 ~
  \href{mailto:arul.rhik@gmail.com}{\raisebox{-0.2\height}\faEnvelope\ arul.rhik@gmail.com} ~
  \href{https://arulrhikm.github.io/}{\raisebox{-0.1\height}{\includegraphics[height=1em]{media/images/github_PNG20.png}} arulrhikm.github.io} ~
  \href{https://www.linkedin.com/in/arul-mazumder-370294239/}{\raisebox{-0.2\height}\faLinkedin\ Arul (Rhik) Mazumder} ~
  \href{https://scholar.google.com/citations?user=NU6USH0AAAAJ&hl=en}{\raisebox{-0.2\height}{\includegraphics[height=1em]{media/images/google-scholar.png}} Arul Mazumder}
\end{center}
\vspace{-6pt}

\small
\textbf{Research interests.} I study what quantum algorithms really cost and when they are worth it: making
algorithms such as Hamiltonian simulation cheaper on fault-tolerant hardware, proving how many T gates and ancilla
qubits any circuit must use, and understanding when quantum approaches can genuinely
outperform classical algorithms.
\normalsize

%----------EDUCATION----------
\section{Education}
\entry{Imperial College London}{Sep. 2026 -- Sep. 2027 (expected)}{MSc, Physics with Quantum Dynamics}{London, UK}
\entry{Carnegie Mellon University}{Aug. 2023 -- May 2026}{B.S. in Computer Science; ML and Algorithms \& Complexity concentrations}{Pittsburgh, PA}
\begin{itemize}\small
  \item Graduated with University Honors. Relevant coursework: Quantum Computing, Complexity Theory, Learning
  Theory, Randomized Algorithms, Convex Optimization, Parallel Computer Architecture.
\end{itemize}

%----------RESEARCH----------
\section{Research Experience}
\entry{Institute for Quantum Computing (IQC), University of Waterloo}{May 2026 -- Aug. 2026}{Undergraduate Researcher (Mitacs Globalink Scholar), hosted by Michele Mosca and Vlad Gheorghiu}{Waterloo, ON, Canada}
\begin{itemize}\small
  \item Developed tighter T-count bounds for structured unitaries: polynomial-time lower bounds for diagonal
  Clifford+T gates, computed from the gate's phase polynomial and exact for layers of disjoint Toffoli-type gates
  ($6m+1$ T gates) in the $\{\mathrm{CNOT},T\}$ model with clean ancillas. Single-author manuscript in preparation; presented at QRE 2026 and UnitaryCon 2026.
  \item Paired the floor with an $O(n^3)$ upper bound that brackets every gate; the bounds carry over to Toffoli
  networks and finite-field multiplication oracles, certifying a published $\mathrm{GF}(4)$ multiplier optimal at
  17 T gates.
  \item Completed USEQIP: quantum information coursework plus 30+ hours of hands-on experiments on IQC's quantum hardware.
\end{itemize}

\entry{Caltech Institute for Quantum Information and Matter (IQIM)}{May 2025 -- Aug. 2025}{SURF Undergraduate Researcher, Preskill group, with Samson Wang and James Watson (Google Quantum AI)}{Pasadena, CA}
\begin{itemize}\small
  \item Worked on early fault-tolerant algorithms for matrix functions, using commutator scaling to lower the cost
  of product-formula routines; written up as \href{https://arxiv.org/abs/2608.13862}{arXiv:2608.13862} and presented as a poster
  at TQC 2026.
  \item Joined a biweekly reading group on quantum signal processing (QSP) and quantum singular value transformation (QSVT).
\end{itemize}

\entry{BlueQubit}{Oct. 2025 -- Present}{Part-time Researcher and Hackathon Lead, AI and Core teams}{Los Angeles, CA (remote)}
\begin{itemize}\small
  \item Benchmarked GPU-accelerated MPS and Pauli-path simulators against AWS Braket, Quantum Rings and
  open-source backends; the runtime scaling study was released as a
  \href{https://arxiv.org/abs/2607.09882}{preprint} and informed product decisions and
  simulator cost estimators.
  \item Lead student and community engagement: turn BlueQubit research into accessible, technically accurate blog
  posts and tutorials, and mentor students and judge the BlueQubit track at student quantum hackathons.
\end{itemize}

\entry{Quantum Technologies Group, Carnegie Mellon University}{May 2023 -- Jun. 2025}{Researcher with Sridhar Tayur}{Pittsburgh, PA}
\begin{itemize}\small
  \item Wrote a tutorial paper with five accompanying
  \href{https://github.com/arulrhikm/Solving-QUBOs-on-Quantum-Computers}{notebooks} on posing and solving QUBO
  problems with QAOA and quantum annealing (\textit{Tutorials in Operations Research}, INFORMS 2025, pp.~145--183).
  \item Implemented and tested semidefinite-programming formulations for entanglement detection using CVXPY,
  PICOS and Gurobi.
\end{itemize}

\entry{Algorithms Group, George Mason University}{May 2022 -- Jun. 2025}{Researcher with Fei Li}{Remote}
\begin{itemize}\small
  \item Implemented quantum string matching in Qiskit, benchmarked against Rabin--Karp and Knuth--Morris--Pratt (ACM AISS 2022).
  \item Built a hybrid quantum-classical algorithm for the traveling salesman problem on D-Wave QPUs that outperformed
  the pure classical and pure quantum approaches it was benchmarked against (IEEE IPDPS 2023).
  \item Designed three classical heuristics and a D-Wave hybrid algorithm for shipment rerouting, tested with up to
  1,000 shipments on real transport networks (Chicago, Barcelona, Winnipeg) against classical solvers (JSSR 2024).
\end{itemize}

%----------PUBLICATIONS----------
\section{Publications}
\small
\subsection*{Manuscripts and preprints}
\begin{enumerate}
  \item \textbf{A. R. Mazumder}, J. Watson, S. Wang. ``\href{https://arxiv.org/abs/2608.13862}{Resource-Efficient
  Quantum Eigenvalue Transform with Commutator Scaling}.'' arXiv:2608.13862; \textit{Theory of Quantum Computation,
  Communication and Cryptography (TQC)}, 2026, poster.
  \item \textbf{A. R. Mazumder}. ``The Cubic Moment Form.'' Manuscript in preparation, 2026. T-count lower bounds
  for diagonal gates.
  \item J. Lee, \textbf{A. R. Mazumder}. ``How Much Workspace Does Optimal-T Synthesis Need? Ancilla Complexity in
  the Mixed Clifford+T Model.'' Manuscript in preparation, 2026. Ancilla bounds for optimal-T synthesis.
  \item \textbf{A. R. Mazumder}, H. Karapetyan, M. Z. Mullath, R. H. Tanin, H. Gharibyan, H. Tepanyan.
  ``\href{https://arxiv.org/abs/2607.09882}{Benchmarking Zero-Setup Quantum Circuit
  Simulators}.'' arXiv:2607.09882, 2026. Runtime scaling study of GPU-accelerated circuit simulators.
  \item F. Li, \textbf{A. R. Mazumder}, M. Zhao. ``\href{https://arxiv.org/pdf/2501.05624}{Quantum Annealing
  Approaches to Shipment Rerouting Problems}.'' arXiv:2501.05624, 2025.
\end{enumerate}

\subsection*{Peer-reviewed}
\begin{enumerate}
  \item \textbf{A. R. Mazumder}. ``\href{https://arulrhikm.github.io/media/documents/qdgmbench.pdf}{An End-to-End
  Hybrid Quantum--Classical Sampling Workflow for Discrete Markov Random Fields}.'' \textit{IEEE International Conference on Quantum Computing and Engineering (QCE)}, 2026.
  \item \textbf{A. R. Mazumder}, S. R. Mazumder. ``\href{https://arulrhikm.github.io/media/documents/qtda.pdf}{Depth-Efficient
  Quantum Topological Data Analysis for Regime-Specific Detection of Financial Stress}.'' \textit{IEEE International Conference on Quantum Computing and Engineering (QCE)}, 2026.
  \item \textbf{A. R. Mazumder}, S. Tayur. ``\href{https://arxiv.org/pdf/2401.08989}{Five Starter Problems: Solving
  Quadratic Unconstrained Binary Optimization Models on Quantum Computers}.'' \textit{Tutorials in Operations
  Research}, INFORMS, 2025, pp.~145--183.
  \item M. Zhao, \textbf{A. R. Mazumder}, F. Li. ``\href{https://doi.org/10.13021/jssr2024.4263}{Improved Classical
  and Quantum Algorithms for the Shipment Rerouting Problems}.'' \textit{Journal of Student-Scientists' Research}, 2024, doi:10.13021/jssr2024.4263.
  \item \textbf{A. R. Mazumder}, A. Sen, U. Sen. ``\href{https://arxiv.org/pdf/2309.16796v3}{Benchmarking
  Metaheuristic-Integrated QAOA Against Quantum Annealing}.'' \textit{LNNS}, 2024.
  \item F. Li, \textbf{A. R. Mazumder}. ``\href{https://ieeexplore.ieee.org/abstract/document/10177447}{An Adaptive
  Hybrid Quantum Algorithm for the Metric Traveling Salesman Problem}.'' \textit{IPDPS} 2023.
  \item M. Gao, R. Huang, \textbf{A. R. Mazumder}, F. Li. ``\href{https://dl.acm.org/doi/pdf/10.1145/3573834.3574498}{Comparisons
  of Classic and Quantum String Matching Algorithms}.'' \textit{AISS} 2022.
\end{enumerate}
\normalsize

%----------TALKS----------
\section{Talks and Posters}
\begin{itemize}\small
  \item \textbf{UnitaryCon 2026}, Toronto --- poster, ``Certified T-count floors for Toffoli networks and reversible circuits.''
  \item \textbf{Quantum Resource Estimation (QRE) 2026}, Toronto --- poster, same title.
  \item \textbf{IEEE QCE 2026}, Toronto --- two first-author papers.
  \item \textbf{TQC 2026} --- poster, ``Resource-Efficient Quantum Eigenvalue Transform with Commutator Scaling''
  (\href{https://arulrhikm.github.io/media/documents/tqc2026-commutator-scaling-poster.pdf}{PDF}).
  \item \textbf{Quantum Leap Career Nexus}, University of Maryland, Nov. 2025, and \textbf{Caltech IQIM SURF Seminar},
  Aug. 2025 --- talks, ``Early Fault-Tolerant Quantum Algorithms for Matrix Functions via Trotter Extrapolation''
  (\href{https://arulrhikm.github.io/media/documents/surf2025-matrix-functions-slides.pdf}{slides}).
  \item \textbf{INFORMS Annual Meeting 2025} --- \href{https://arulrhikm.github.io/media/slides/QUBOonQCslides.pdf}{presentation}
  on solving QUBO models on quantum computers.
\end{itemize}

%----------AWARDS----------
\section{Awards and Honors}
\begin{itemize}\small
  \item \textbf{IEEE QCE Student Travel Grant} (2026), IEEE Quantum Technical Community.
  \item \textbf{Mitacs Globalink Research Award and IQC Undergraduate Research Award} (2026), funding the IQC
  research term.
  \item \textbf{Duke STAQ NSF-funded Invited Scholar} (2025): week-long PhD-level seminar on quantum simulation,
  error correction and architecture with Peter Love, Kenneth Brown and Yu Tong, including a Duke Quantum Center
  lab tour.
  \item \textbf{MIT iQuHACK} (2026): 2nd, NVIDIA track, of 90+ teams; variational LABS with Pauli correlation encodings on GPUs.
  \item \textbf{YQuantum} (2025, Yale): Yale Quantum Institute Grand Prize; 1st place, Alice \& Bob track;
  2nd overall; Wigner-function quantum state reconstruction via convex optimization and convolutional-network
  denoising.
  \item \textbf{Citadel/Correlation One Datathon} (2024): 2nd of 150 invited from 3,000+ applicants; pairs trading, Sharpe 1.58.
  \item \textbf{HackCMU} (2023): Graduate Student Association Infrastructure Prize for Eco-Bin, an AI waste sorter
  piloted with CMU Dining.
  \item \textbf{Regeneron Science Talent Search Scholar} (2023), top 300 of 2,000+ applicants nationally.
  \item \textbf{Olympiads:} USA(J)MO qualifier, USAPhO Bronze, USACO Gold, USNCO National Finalist.
\end{itemize}

%----------LEADERSHIP----------
\section{Leadership}
\entry{CMU Quantum}{Sep. 2024 -- May 2026}{Founder and President}{Pittsburgh, PA}
\begin{itemize}\small
  \item Founded CMU's first student quantum organization and grew it to 200+ members and 10+ student leaders;
  secured sponsorship from General Motors and Boeing; ran workshops with IBM, IonQ and Quantinuum.
  \item Developed teaching materials (\href{https://drive.google.com/file/d/1V5-msJkDjIr6-rSO2Juw5q21DpHwTdrh/view?usp=sharing}{slides and handouts}),
  led hackathon tracks, and organized sponsored research trips.
\end{itemize}

%----------SKILLS----------
\section{Technical Skills}
\small
\textbf{Languages and proof tools:} Python, Julia, C++, Lean 4\\
\textbf{Quantum and optimization:} Qiskit, PennyLane, CUDA-Q, D-Wave Ocean; CVXPY, PICOS, Gurobi; CUDA, MPI
\normalsize

\end{document}
